\subsection{给家长：分龄谈性与就医陪伴的要点}

前面的问答是给青少年自己看的；这一小节说给家长。研究一再显示：家长若能\textbf{平静、准确、适龄}地与孩子谈性，孩子反而会更晚发生首次性行为、更懂得保护自己——"谈性"不会教坏孩子，沉默才会把孩子推向不靠谱的信息源。

\begin{itemize}
	\item \textbf{小学阶段（约 6--12 岁）}：回答"我从哪里来"用科学而简短的事实（卵子、精子、子宫），不编"垃圾桶捡来的"；教会孩子\textbf{科学部位名称}与"泳衣遮盖的地方不许别人碰"，这是防性侵教育最有效的年龄窗口
	\item \textbf{初中阶段（约 12--15 岁）}：主动谈青春期的身体变化（遗精、月经、自慰）与情感萌动；重要的是\textbf{先听后讲}——"你同学们都聊些什么？"比一上来就训诫更能打开话匣子
	\item \textbf{高中阶段（约 15--18 岁）}：明确而务实地谈避孕、安全套与同意（consent），并告诉孩子"如果出了状况，来找我，我们一起处理"。研究显示，被许诺"求助不会先挨骂"的青少年，\textbf{更可能在第一时间求助}，而不是隐瞒拖延
	\item \textbf{就医陪伴}：带青春期孩子就诊（月经紊乱、外阴炎、包皮问题、性发育疑问）时，尊重孩子的隐私意愿——多数医院允许青少年与医生单独交流几分钟；家长在场时不要抢答或斥责
	\item \textbf{发现"早恋"或性活动迹象}：先稳住自己的情绪。沟通目标是\textbf{安全与知情}（是否自愿、是否防护、年龄是否合法），而不是羞辱与没收手机——羞辱只会把孩子推向更深的隐瞒
\end{itemize}

\begin{tcolorbox}[colback=pink!5!white,colframe=red!50!black,title=贴心小叮咛]
家长不需要"什么都会答"。答不上来时，诚实地说"这个我查一下再告诉你"，再一起查权威资料——\textbf{和孩子一起学习性知识，本身就是最好的性教育}。
\end{tcolorbox}

